% Generated draft: requires mathematical and source review.
\documentclass[11pt]{article}
\usepackage[margin=25mm]{geometry}
\usepackage{fontspec}
\setmainfont{Latin Modern Roman}
\usepackage{amsmath,amssymb,mathtools,unicode-math}
\setmathfont{Latin Modern Math}
\usepackage{enumitem,needspace,xurl,hyperref}
\setlength{\parindent}{0pt}
\setlength{\parskip}{5pt}
\setlength{\emergencystretch}{3em}
\newcommand{\N}{\mathbb{N}}
\newcommand{\Z}{\mathbb{Z}}
\newcommand{\Q}{\mathbb{Q}}
\newcommand{\R}{\mathbb{R}}
\newcommand{\C}{\mathbb{C}}
\newcommand{\F}{\mathbb{F}}
\newcommand{\A}{\mathbb{A}}
\newcommand{\PP}{\mathbb{P}}
\newcommand{\E}{\mathbb{E}}
\newcommand{\eps}{\varepsilon}
\newcommand{\dd}{\,\mathrm{d}}
\begin{document}
\section*{UnsolvedMath ID: 30000429}
\textbf{Source problem code:} catalogue entry
\subsection{Definitions and mathematical statement}
For a discrete group \(G\), let \(\C G\) be its complex group algebra. Let \(\mathcal N(G)\) be the group von Neumann algebra in the left regular representation on \(\ell^2(G)\), and let \(U(G)\) denote the ring of operators affiliated with \(\mathcal N(G)\), with its standard algebraic operations. The source asks whether there is a unital division subring \(D\subseteq U(G)\) such that
\[
\C G\subseteq D\subseteq U(G).
\]
A skew field means that every nonzero element has a two-sided multiplicative inverse; commutativity is not required. The inclusions use the canonical group-algebra representation and are not required to be proper. An abstract embedding into a division ring unrelated to \(U(G)\) would be a different target.

The source gives no hypothesis on \(G\). Torsion-freeness is a necessary qualification: if \(g\in G\) has finite order \(m>1\), then the two nonzero group-algebra elements satisfy
\[
(1-g)(1+g+\cdots+g^{m-1})=0,
\]
which is incompatible with an embedding into a skew field. Thus a universal version over arbitrary groups cannot be the intended positive conjecture; whether \(G\) is assumed torsion-free or belongs to a more restricted class must be recovered rather than silently supplied.

Let \(\mathcal D_G\) denote the division closure of \(\C G\) in \(U(G)\), namely the smallest subring containing \(\C G\) and closed under inversion of its elements that are invertible in \(U(G)\). On the given embedding, the existence question can be expressed as whether \(\mathcal D_G\) is a skew field. This formulation does not assume that every nonzero element of an arbitrary affiliated-operator subring is invertible.

The supplied assessment checked 2026-08-21 reports that division closures are skew fields for important torsion-free group classes, but not settlement of the unrestricted extracted existence question. It does not identify the exact classes or provide a theorem for all torsion-free groups. The literal missing torsion assumption and the unresolved intended scope are therefore recorded separately from that dated status.
\subsection{Short English statement}
Can the canonical complex group algebra sit inside a division ring contained in its affiliated-operator ring? The missing hypothesis on the group is essential: torsion rules out such an embedding, while the supplied review concerns partial results for torsion-free classes.
\subsection{Sources}
\begin{itemize}
\item Local UnsolvedMath registry, permanent ID 30000429, catalogue entry, \emph{Skew Fields between \(\C G\) and \(U(G)\)}.
\item \emph{Miniworkshop: L2-Spectral Invariants and the Integrated Density of States} (2006), Oberwolfach Reports, pp. 511--552, DOI \url{https://doi.org/10.4171/owr/2006/09}.
\item Supplied reference ``Linnell--L\"uck localization'': \url{https://arxiv.org/abs/1602.06906}.
\item Supplied literature triage checked 2026-08-21; no exact torsion-free group-class hypotheses or separate theorem statement is reproduced.
\end{itemize}
\end{document}
